\chapter{Fanno Flow}

Fanno flow describes compressible, adiabatic flow through a constant-area duct where friction is the dominant effect. Friction continuously drives the flow toward the sonic condition (Mach 1), the critical point. For subsonic inlet flow, friction accelerates the flow toward $M = 1$; for supersonic inlet flow, friction decelerates the flow toward $M = 1$. The friction parameter $fL^*/D$ (where $f$ is the Darcy friction factor, $L^*$ is the length needed to reach sonic conditions, and $D$ is the duct diameter) quantifies how much duct is available relative to the critical length. Fanno flow is widely used in modelling pipe flow in propulsion systems and high-speed internal flows. This module was recreated from the legacy toolkit codebase originally developed by Gowtham Sivaraman \cite{sivaraman2015}.

\section{Nomenclature}

\begin{center}
\begin{tabular}{@{}cl@{}}
\toprule
Symbol & Description \\
\midrule
$\gamma$ & Specific heat ratio \\
$M$ & Mach number \\
$\frac{T}{T^*}$ & Static to Sonic Temperature ratio \\
$\frac{P}{P^*}$ & Static to Sonic Pressure ratio \\
$\frac{\rho}{\rho^*}$ & Density to sonic density ratio \\
$\frac{P_0}{P_0^*}$ & Local stagnation to Sonic stagnation Pressure ratio \\
$\frac{f L^*}{D}$ & Friction parameter (limiting duct length to reach sonic conditions) \\
$\frac{U}{U^*}$ & Velocity to sonic velocity ratio \\
\bottomrule
\end{tabular}
\end{center}

\section{Governing Idea}

Every Fanno-flow property is a function of specific heat ratio $\gamma$ and Mach number $M$ alone. The solver therefore treats $M$ as the pivot variable: whichever quantity the user supplies is first inverted to recover $M$, after which every remaining property is evaluated directly from $M$ and $\gamma$. The problem is split into seven cases depending on which property is given.

\section{Calculation Cases}

Depending on the input parameter provided by the user, the Fanno flow properties are computed using one of the seven cases described below.

\subsection*{Case 1: Given $\gamma$ and $M$}

\textbf{Given:}
\begin{itemize}[nosep]
    \item Specific heat ratio $\gamma > 1$
    \item Mach number $M > 0$
\end{itemize}

\textbf{Calculation:} All properties follow directly from explicit analytical relations:
\begin{equation}
\frac{T}{T^*} = \frac{\frac{\gamma+1}{2}}{1+\frac{\gamma-1}{2}M^2}
\end{equation}

\begin{equation}
\frac{P}{P^*} = \frac{1}{M}\sqrt{\frac{\frac{\gamma+1}{2}}{1+\frac{\gamma-1}{2}M^2}}
\end{equation}

\begin{equation}
\frac{\rho}{\rho^*} = \frac{1}{M}\sqrt{\frac{1+\frac{\gamma-1}{2}M^2}{\frac{\gamma+1}{2}}}
\end{equation}

\begin{equation}
\frac{P_0}{P_0^*} = \frac{1}{M}\left(\frac{1+\frac{\gamma-1}{2}M^2}{\frac{\gamma+1}{2}}\right)^{\frac{\gamma+1}{2(\gamma-1)}}
\end{equation}

\begin{equation}
\frac{f L^*}{D} = \frac{\gamma+1}{2\gamma}\ln\left[\frac{\frac{\gamma+1}{2}M^2}{1+\frac{\gamma-1}{2}M^2}\right] + \frac{1}{\gamma}\left(\frac{1}{M^2}-1\right)
\end{equation}

\begin{equation}
\frac{U}{U^*} = M\sqrt{\frac{\frac{\gamma+1}{2}}{1+\frac{\gamma-1}{2}M^2}}
\end{equation}


\subsection*{Case 2: Given $\gamma$ and $\frac{T}{T^*}$}

\textbf{Given:}
\begin{itemize}[nosep]
    \item Specific heat ratio $\gamma > 1$
    \item Temperature ratio $0 < \frac{T}{T^*} < \left(\frac{T}{T^*}\right)_{max}$, where the maximum temperature ratio is:
    \begin{equation}
    \left(\frac{T}{T^*}\right)_{max} = \frac{\gamma+1}{2}
    \end{equation}
\end{itemize}

\textbf{Calculation:} The Mach number $M$ is solved explicitly in closed form:
\begin{equation}
M = \sqrt{\frac{\gamma+1-2\frac{T}{T^*}}{\frac{T}{T^*}(\gamma-1)}}
\end{equation}


\subsection*{Case 3: Given $\gamma$ and $\frac{P}{P^*}$}

\textbf{Given:}
\begin{itemize}[nosep]
    \item Specific heat ratio $\gamma > 1$
    \item Pressure ratio $\frac{P}{P^*} > 0$
\end{itemize}

\textbf{Calculation:} The Mach number is solved explicitly in closed form:
\begin{equation}
M = \sqrt{\frac{\sqrt{\left(\frac{P}{P^*}\right)^2+\gamma^2-1}-\frac{P}{P^*}}{\frac{P}{P^*}(\gamma-1)}}
\end{equation}


\subsection*{Case 4: Given $\gamma$ and $\frac{\rho}{\rho^*}$}

\textbf{Given:}
\begin{itemize}[nosep]
    \item Specific heat ratio $\gamma > 1$
    \item Density ratio $\frac{\rho}{\rho^*} > \left(\frac{\rho}{\rho^*}\right)_{min}$, where the minimum density ratio is:
    \begin{equation}
    \left(\frac{\rho}{\rho^*}\right)_{min} = \sqrt{\frac{\gamma-1}{\gamma+1}}
    \end{equation}
\end{itemize}

\textbf{Calculation:} The Mach number is solved explicitly in closed form:
\begin{equation}
M = \sqrt{\frac{1}{\left(\frac{\rho}{\rho^*}\right)^2\times\frac{\gamma+1}{2}-\frac{\gamma-1}{2}}}
\end{equation}


\subsection*{Case 5: Given $\gamma$ and $\frac{P_0}{P_0^*}$}

\textbf{Given:}
\begin{itemize}[nosep]
    \item Specific heat ratio $\gamma > 1$
    \item Stagnation pressure ratio $\frac{P_0}{P_0^*} \ge 1$
\end{itemize}

\textbf{Calculation:} No closed form exists; the Mach number $M$ is solved iteratively using the bisection method based on the selected flow branch toggle:
\begin{equation}
M = \frac{M_{start}+M_{end}}{2}
\end{equation}
\begin{equation}
f(M) = \frac{1}{M}\left(\frac{1+\frac{\gamma-1}{2}M^2}{\frac{\gamma+1}{2}}\right)^{\frac{\gamma+1}{2(\gamma-1)}} - \frac{P_0}{P_0^*} = 0
\end{equation}
The bisection solver updates the search brackets until $|f(M)| < 10^{-9}$ (with a limit of 100 iterations), using the following brackets based on the active branch:
\begin{itemize}[nosep]
    \item \textbf{Subsonic branch toggle:} Search interval $[M_{start}, M_{end}] = [10^{-8}, 1.0]$.
    \item \textbf{Supersonic branch toggle:} Search interval $[M_{start}, M_{end}] = [1.0, 100.0]$.
\end{itemize}


\subsection*{Case 6: Given $\gamma$ and $\frac{f L^*}{D}$}

\textbf{Given:}
\begin{itemize}[nosep]
    \item Specific heat ratio $\gamma > 1$
    \item Friction parameter $\frac{f L^*}{D} > 0$. If the supersonic branch toggle is active, the friction parameter must be less than the theoretical supersonic maximum limit:
    \begin{equation}
    \left(\frac{f L^*}{D}\right)_{sup,max} = \frac{\gamma+1}{2\gamma}\ln\left[\frac{\gamma+1}{\gamma-1}\right] - \frac{1}{\gamma}
    \end{equation}
\end{itemize}

\textbf{Calculation:} No closed form exists; the Mach number $M$ is solved iteratively using the bisection method based on the selected flow branch toggle:
\begin{equation}
M = \frac{M_{start}+M_{end}}{2}
\end{equation}
\begin{equation}
f(M) = \frac{\gamma+1}{2\gamma}\ln\left[\frac{\frac{\gamma+1}{2}M^2}{1+\frac{\gamma-1}{2}M^2}\right] + \frac{1}{\gamma}\left(\frac{1}{M^2}-1\right) - \frac{f L^*}{D} = 0
\end{equation}
The bisection solver updates the search brackets until $|f(M)| < 10^{-9}$ (with a limit of 100 iterations), using the following brackets based on the active branch:
\begin{itemize}[nosep]
    \item \textbf{Subsonic branch toggle:} Search interval $[M_{start}, M_{end}] = [10^{-8}, 1.0]$.
    \item \textbf{Supersonic branch toggle:} Search interval $[M_{start}, M_{end}] = [1.0, 100.0]$.
\end{itemize}


\subsection*{Case 7: Given $\gamma$ and $\frac{U}{U^*}$}

\textbf{Given:}
\begin{itemize}[nosep]
    \item Specific heat ratio $\gamma > 1$
    \item Velocity ratio $0 < \frac{U}{U^*} < \left(\frac{U}{U^*}\right)_{max}$, where the maximum velocity ratio is:
    \begin{equation}
    \left(\frac{U}{U^*}\right)_{max} = \sqrt{\frac{\gamma+1}{\gamma-1}}
    \end{equation}
\end{itemize}

\textbf{Calculation:} The Mach number is solved explicitly in closed form:
\begin{equation}
M = \frac{2\times\frac{U}{U^*}}{\sqrt{2\gamma+2-\left(2\left(\frac{U}{U^*}\right)^2\gamma\right)+2\left(\frac{U}{U^*}\right)^2}}
\end{equation}



\section{Program Flow and Architecture}

In Fanno flow, $M$ is the pivot variable. Five inputs use explicit inversion; $P_0/P_0^*$ and $fL^*/D$ use bisection with the sub/supersonic toggle.

\begin{center}
\begin{tikzpicture}[x=1cm,y=1cm,
    inp/.style={rectangle, draw=blue!70!black, thick, fill=blue!6, align=center,
                text width=3.8cm, minimum height=0.65cm, font=\small\sffamily},
    eng/.style={rectangle, draw=orange!80!black, thick, fill=orange!8, align=center,
                text width=3.4cm, font=\small\sffamily},
    outn/.style={rectangle, draw=green!60!black, thick, fill=green!6, align=center,
                text width=2.4cm, minimum height=0.65cm, font=\small\sffamily},
    note/.style={rectangle, draw=gray!60, dashed, fill=gray!5, align=center,
                 text width=2.0cm, minimum height=0.55cm, font=\scriptsize\sffamily},
    arr/.style={-{Stealth[scale=1.0]}, thick, draw=black!70},
    darr/.style={-{Stealth[scale=1.0]}, thick, dashed, draw=black!50},
    bus/.style={thick, draw=orange!70!black}
]
\node[inp] (i0) at (0, 0.0) {Case 1: $M$ (direct)};
\node[inp] (i1) at (0, 0.9) {Case 2: $T/T^*$ $\to$ explicit};
\node[inp] (i2) at (0, 1.8) {Case 3: $P/P^*$ $\to$ explicit};
\node[inp] (i3) at (0, 2.7) {Case 4: $\rho/\rho^*$ $\to$ explicit};
\node[inp] (i4) at (0, 3.6) {Case 5: $P_0/P_0^*$ $\to$ bisection};
\node[inp] (i5) at (0, 4.5) {Case 6: $fL^*/D$ $\to$ bisection};
\node[inp] (i6) at (0, 5.4) {Case 7: $U/U^*$ $\to$ explicit};
\node[note] (t1) at (4.3, 3.6) {Sub/Sup\\toggle};
\node[note] (t2) at (4.3, 4.5) {Sub/Sup\\toggle};
\node[eng, minimum height=6.0cm] (eng) at (7.2, 2.7)
    {\textbf{FannoFlowEngine}\\[3pt]
    \texttt{fromMach(M,\,$\gamma$)}};
\node[outn] (o0) at (11.8, 0.0) {$M$};
\node[outn] (o1) at (11.8, 0.9) {$T/T^*$};
\node[outn] (o2) at (11.8, 1.8) {$P/P^*$};
\node[outn] (o3) at (11.8, 2.7) {$\rho/\rho^*$};
\node[outn] (o4) at (11.8, 3.6) {$P_0/P_0^*$};
\node[outn] (o5) at (11.8, 4.5) {$fL^*/D$};
\node[outn] (o6) at (11.8, 5.4) {$U/U^*$};
\foreach \i in {i0,i1,i2,i3,i4,i5,i6} {\draw[arr] (\i.east) -- (eng.west |- \i);}
\draw[darr] (t1.east) -- (i4.east |- t1) -- (i4.east);
\draw[darr] (t2.east) -- (i5.east |- t2) -- (i5.east);
\coordinate (bx) at ($(eng.east)+(0.5,0)$);
\draw[bus] (eng.east) -- (bx);
\draw[bus] (bx |- o0) -- (bx |- o6);
\foreach \n in {o0,o1,o2,o3,o4,o5,o6} {\draw[arr] (bx |- \n) -- (\n.west);}
\end{tikzpicture}
\end{center}